\documentclass[a4paper]{article}
\usepackage{graphicx}
\usepackage{amsmath,amssymb}
\usepackage{minted}
\usepackage{multirow}
\usepackage{float}
\usepackage{algorithm}
\usepackage{algpseudocode}
\usepackage{todonotes}
\usepackage{forest}
\usepackage{xstring}
\usepackage{xcolor}
\usepackage[most]{tcolorbox}
\usepackage{xparse}
\usepackage{natbib}
\usetikzlibrary{graphs,graphdrawing}
\usegdlibrary{layered}
% ND Mind Mirror default input files.
\input{/home/donkarlo/Dropbox/repo/nd_language_ai_project/src/nd_language_ai/formal/computer/programming/tex/latex/code_clips/external_graphviz.tex}
\input{/home/donkarlo/Dropbox/repo/nd_language_ai_project/src/nd_language_ai/formal/computer/programming/tex/latex/part/control_sequence/kind/semtag/semtag.tex}
\input{/home/donkarlo/Dropbox/repo/nd_language_ai_project/src/nd_language_ai/formal/computer/programming/tex/latex/code_clips/seehere.tex}

\title{LOESS and LOWESS Trend Smoothing}
\date{}
\author{Mohammad Rahmani}
\begin{document}
\maketitle

\section{Definition}
LOWESS (locally weighted scatterplot smoothing) and LOESS (locally estimated scatterplot smoothing) fit simple polynomial models in a neighborhood of each target point. The terminology overlaps, while LOESS often covers a broader choice of local polynomial degrees.

\section{Methodology}
At target time $t_0$, use a fraction of the nearby observations with distance-based weights. A common tricube weight is
\[
 w_i(t_0)=
 \begin{cases}
 \left(1-\left|\frac{t_i-t_0}{h}\right|^3\right)^3,
       &|t_i-t_0|<h,\\
 0,&\text{otherwise},
 \end{cases}
\]
where $h$ describes the neighborhood radius.

Fit a local linear model by weighted least squares:
\[
 (\widehat a,\widehat b)=\underset{a,b}{\arg\min}\,
 \sum_i w_i(t_0)\bigl[y_i-a-b(t_i-t_0)\bigr]^2.
\]
The local smoothed value is $\widehat T(t_0)=\widehat a$. One can instead fit a local quadratic polynomial. Some implementations also add residual-based robustness iterations.

\section{E-Commerce Example}
Demand for a new product might grow quickly after its launch and gradually flatten. A straight line would miss the changing growth rate, whereas a suitably tuned LOESS smoother can display the bend.

\section{Python Illustration}
\begin{verbatim}
from statsmodels.nonparametric.smoothers_lowess import lowess

trend = lowess(
    endog=revenue,
    exog=day_number,
    frac=0.20,
    it=2,
    return_sorted=False,
)
\end{verbatim}
The fraction controls the neighborhood size: larger values give smoother curves, but can blur short-term turning points.

\section{Advantages}
\begin{itemize}
 \item Reveals nonlinear smooth relationships without selecting a global function.
 \item Flexible for exploratory trend visualization.
 \item Supports local rather than one-size-fits-all fitting.
\end{itemize}

\section{Limitations}
\begin{itemize}
 \item Sensitive to span/fraction and polynomial degree.
 \item Endpoints are less stable and standard fits may use future observations.
 \item Does not explicitly isolate seasonality or establish trend significance.
 \item Repeated local fits may be expensive for very large datasets.
\end{itemize}

\section{Practical Guidance}
Plot the raw values along with the fitted curve and check sensitivity to multiple spans. Remove or model strong seasonal effects first. For prediction-time use, avoid future data by using historical-only fitting or explicitly one-sided neighborhoods.

\bibliographystyle{plainnat}
\bibliography{/home/donkarlo/Dropbox/repo/nd_shared_research/bibliography/bibliography.bib}
\end{document}
